% Generated draft: requires mathematical and source review.
\documentclass[11pt]{article}
\usepackage[margin=25mm]{geometry}
\usepackage{fontspec}
\setmainfont{Latin Modern Roman}
\usepackage{amsmath,amssymb,mathtools,unicode-math}
\setmathfont{Latin Modern Math}
\usepackage{enumitem,needspace,xurl,hyperref}
\setlength{\parindent}{0pt}
\setlength{\parskip}{5pt}
\setlength{\emergencystretch}{3em}
\newcommand{\N}{\mathbb{N}}
\newcommand{\Z}{\mathbb{Z}}
\newcommand{\Q}{\mathbb{Q}}
\newcommand{\R}{\mathbb{R}}
\newcommand{\C}{\mathbb{C}}
\newcommand{\F}{\mathbb{F}}
\newcommand{\A}{\mathbb{A}}
\newcommand{\PP}{\mathbb{P}}
\newcommand{\E}{\mathbb{E}}
\newcommand{\eps}{\varepsilon}
\newcommand{\dd}{\,\mathrm{d}}
\begin{document}
\section*{UnsolvedMath ID: 30005094}
\textbf{Source problem code:} catalogue entry
\subsection{Definitions and mathematical statement}
The supplied record concerns the following conjecture, with two conventions requiring confirmation: it does not explicitly define matrix height or state the full range of matrix dimensions. Under the usual entrywise-height interpretation, for a fixed positive integer \(n\) and a real number \(H\geq 1\), put
\[
\mathcal M_n(\Z;H)=\{A=(a_{ij})\in\operatorname{Mat}_{n\times n}(\Z):\max_{1\leq i,j\leq n}|a_{ij}|\leq H\}.
\]
For a monic polynomial \(f\in\Z[X]\) of degree \(n\), define
\[
R_n(H;f)=\#\{A\in\mathcal M_n(\Z;H):\det(XI_n-A)=f(X)\},
\]
where \(I_n\) is the identity matrix. Restricting attention to monic integral polynomials of degree \(n\) retains every possible characteristic polynomial of an integer \(n\times n\) matrix.

\emph{Conjecture 1 in the supplied excerpt.} For each fixed dimension \(n\) in the intended range,
\[
R_n(H;f)\leq H^{n(n-1)/2+o(1)}\qquad(H\to\infty),
\]
uniformly over \(f\). Explicitly, uniformity means that for every \(\varepsilon>0\) there exist constants \(C_{n,\varepsilon}>0\) and \(H_{n,\varepsilon}\geq1\), independent of \(f\), such that
\[
\forall H\geq H_{n,\varepsilon}\quad
\forall f\in\Z[X]\text{ monic of degree }n,\qquad
R_n(H;f)\leq C_{n,\varepsilon}H^{n(n-1)/2+\varepsilon}.
\]
In particular, \(f\) may vary with \(H\); no uniformity in \(n\) is asserted. The entrywise-height convention and the intended lower bound on \(n\) remain subject to source confirmation.

\emph{Supplied contextual claims, not independently verified here.} The excerpt states that the conjecture holds for \(n=2\), and gives the unconditional uniform estimate
\[
R_n(H;f)\leq H^{n^2-n+o(1)}.
\]
It also introduces the conditional hypothesis that, for some \(\gamma_n\geq0\), uniformly over \(f\),
\[
R_n(H;f)\leq H^{n^2-n-\gamma_n+o(1)}.
\]
The excerpt identifies \(\gamma_n=n(n-1)/2\) with the conjectured exponent, states that \(\gamma_n=0\) is available unconditionally, and reports an explicit positive sequence with \(\gamma_3=\gamma_4=1\).

For integers \(d,t\), the excerpt further considers
\[
S_n(H;d,t)=\#\{A\in\mathcal M_n(\Z;H):\det A=d,\ \operatorname{tr}A=t\}.
\]
Writing \(f(X)=X^n+c_{n-1}X^{n-1}+\cdots+c_0\), it observes that
\[
R_n(H;f)\leq S_n\bigl(H;(-1)^nc_0,-c_{n-1}\bigr).
\]
Its Theorem 3 is stated for \(n\geq4\), uniformly over \(d,t\), and reports
\[
S_4(H;d,t)\ll H^{11+o(1)},\qquad
S_n(H;d,t)\ll H^{n^2-n-\gamma_n}\quad(n\geq5).
\]
The explicit formula for \(\gamma_n\) accompanying this theorem is corrupted in the supplied extraction and cannot be reconstructed reliably. These contextual bounds are recorded as source claims rather than additional targets or independently established facts.
\subsection{Short English statement}
For each fixed matrix size, count integer matrices whose entries have absolute value at most \(H\) and whose characteristic polynomial equals a prescribed polynomial. The conjecture asks whether this count is at most \(H^{n(n-1)/2+o(1)}\), with the error uniform over the polynomial. The height convention and intended range of matrix sizes require confirmation. The supplied record labels the problem open and reports the case \(n=2\) as proved; these are supplied assessments.
\subsection{Sources}
\begin{itemize}
\item Supplied local UnsolvedMath registry record, permanent ID \(30005094\), title \emph{Counting Integer Matrices with Fixed Characteristic Polynomial}, exact source problem code \texttt{catalogue entry}, in the collection \emph{Oberwolfach Reports Open Problems}. The record gives proposed year 2022 and names no proposer.
\item Supplied primary citation: \emph{Diophantische Approximationen (2023)}, Oberwolfach Reports, pp.\ 1091--1164, DOI 10.4171/owr/2022/21. \url{https://doi.org/10.4171/owr/2022/21}. The background describes Oberwolfach Reports as documenting workshops and their open-problem discussions; the literature-source label is ``Oberwolfach Report 21/2022.'' The relationship between the supplied 2023 bibliographic year and the 2022 report designation requires confirmation.
\item Supplied literature source: Ostafe--Shparlinski, ``integer matrices.'' \url{https://doi.org/10.48550/arXiv.2203.03880}. The record supplies no fuller bibliographic description or explicit attribution of the conjecture to these authors.
\item Literature assessment dated exactly \texttt{2026-08-22}: the supplied curation labels the conjecture open, reports that the primary report proves the case \(n=2\), and reports weaker uniform bounds in higher dimensions. Its listed sources checked are the two links above. This is supplied dated curation, not an independently verified literature assessment; the background expressly requests independent verification of time-sensitive status and cited claims.
\item The registry labels the statement \texttt{corrected\_verified} and records a claimed check against the official Oberwolfach report. That verification claim is supplied metadata and has not been independently confirmed here. The excerpt attributes the basic determinant-counting bound to its reference [8, Theorem 4], and attributes ideas used in the proof of Theorem 3 to Katznelson [3], adjusted to the authors' setting and combined with geometry-of-numbers techniques. The identities of references [8] and [3] are not supplied.
\end{itemize}
\end{document}
